\documentclass{article}
\usepackage[utf8]{inputenc}
\usepackage{amsmath}
\usepackage{geometry}
\geometry{margin=2cm}
\begin{document}

\section*{Questão 130 — ENEM 2024 — Ciências da Natureza}

\subsection*{Interpretação e Estratégia}
Esta questão exige identificar em qual órgão do sistema digestório humano se inicia a reação química entre nitrito de sódio e ácido clorídrico, processo que gera compostos nitrosos mutagênicos, conforme descrito no enunciado.

\textbf{Termos importantes:}
\begin{itemize}
  \item \textit{Mutagênicas}: substâncias capazes de causar mutações no material genético das células.
  \item \textit{Reação química}: transformação entre substâncias, formando novos compostos.
\end{itemize}

\textbf{O que a questão pede:}

Determinar o órgão em que ocorre a reação inicial entre nitrito de sódio e ácido clorídrico, produzindo ácido nitroso que, por sua vez, reage com aminas para formar compostos nitrosos mutagênicos.

\textbf{Estratégia para resolver:}

Relacionar o ambiente ácido da reação química com a localização anatômica da produção de ácido clorídrico no trato digestório, deduzindo o órgão onde esse processo começa.

\subsection*{Conteúdo e Mini-revisão}
\begin{tabular}{|p{4cm}|p{6cm}|p{5cm}|}
\hline
\textbf{Conceito} & \textbf{Ideia-chave} & \textbf{Aplicação} \\ \hline
Ácido clorídrico & Presente no suco gástrico do estômago & Local da reação com nitrito de sódio para formar ácido nitroso \\ \hline
Compostos nitrosos mutagênicos & Formados pela reação do ácido nitroso com aminas & Produzidos inicialmente no estômago \\ \hline
Função dos órgãos digestórios & Cada órgão tem funções específicas & Estômago produz suco gástrico ácido fundamental para a reação \\ \hline
\end{tabular}

\bigskip

\textbf{Fundamentos teóricos:}

\begin{itemize}
  \item O estômago secreta suco gástrico contendo alto teor de ácido clorídrico, criando ambiente ácido.
  \item No ambiente ácido do estômago, nitrito de sódio reage com ácido clorídrico formando ácido nitroso.
  \item O ácido nitroso reage com aminas (compostos nitrogenados), presentes em carnes processadas, originando compostos nitrosos mutagênicos.
  \item Outros órgãos como fígado, intestino, pancreas, ou rins não produzem ácido clorídrico e, portanto, não iniciam essa reação.
\end{itemize}

\subsection*{Resolução e "Pulo do Gato"}

O enunciado informa que:
\begin{itemize}
  \item Nitrito de sódio reage com ácido clorídrico para formar ácido nitroso.
  \item Ácido nitroso reage com aminas para formar compostos nitrosos mutagênicos.
\end{itemize}

Sabemos que o ácido clorídrico é secretado pelo estômago em seu suco gástrico. Assim:

\begin{equation*}
\mathrm{NO_2^- + HCl \rightarrow HNO_2}
\end{equation*}

Esse ácido nitroso, reativo, participa das reações para formar os compostos mutagênicos.

Portanto, \textbf{o local inicial dessas reações é o estômago}.

\textbf{Pulo do gato:} lembrar que \textit{apenas o estômago secreta ácido clorídrico em concentração suficiente para iniciar a reação química descrita}.

\subsection*{Armadilhas das Alternativas}

\begin{itemize}
  \item \textbf{A — Rim}: confundir filtragem e desintoxicação com local de produção de ácido clorídrico.
  \item \textbf{B — Fígado}: pensar que metabolismo hepático significa também reação inicial química descrita.
  \item \textbf{C — Intestino}: associar todas as reações digestivas ao intestino, sem distinguir o local da secreção ácida.
  \item \textbf{D — Pâncreas}: associar a produção de enzimas digestivas e bicarbonato à reação com nitrito de sódio, o que é incorreto.
\end{itemize}

\subsection*{Código da Questão}

CIENCIAS_NATUREZA_BIOQUIMICA_001

\bigskip

\noindent\textbf{Resposta final:} alternativa \textbf{E — Estômago}.

\end{document}